% Options for packages loaded elsewhere
\PassOptionsToPackage{unicode}{hyperref}
\PassOptionsToPackage{hyphens}{url}
\PassOptionsToPackage{dvipsnames,svgnames,x11names}{xcolor}
%
\documentclass[
  11pt,
  a4paper,
]{ltjsarticle}
\usepackage{amsmath,amssymb}
\usepackage{iftex}
\ifPDFTeX
  \usepackage[T1]{fontenc}
  \usepackage[utf8]{inputenc}
  \usepackage{textcomp} % provide euro and other symbols
\else % if luatex or xetex
  \usepackage{unicode-math} % this also loads fontspec
  \defaultfontfeatures{Scale=MatchLowercase}
  \defaultfontfeatures[\rmfamily]{Ligatures=TeX,Scale=1}
\fi
\usepackage{lmodern}
\ifPDFTeX\else
  % xetex/luatex font selection
\fi
% Use upquote if available, for straight quotes in verbatim environments
\IfFileExists{upquote.sty}{\usepackage{upquote}}{}
\IfFileExists{microtype.sty}{% use microtype if available
  \usepackage[]{microtype}
  \UseMicrotypeSet[protrusion]{basicmath} % disable protrusion for tt fonts
}{}
\makeatletter
\@ifundefined{KOMAClassName}{% if non-KOMA class
  \IfFileExists{parskip.sty}{%
    \usepackage{parskip}
  }{% else
    \setlength{\parindent}{0pt}
    \setlength{\parskip}{6pt plus 2pt minus 1pt}}
}{% if KOMA class
  \KOMAoptions{parskip=half}}
\makeatother
\usepackage{xcolor}
\usepackage[margin=24mm]{geometry}
\setlength{\emergencystretch}{3em} % prevent overfull lines
\providecommand{\tightlist}{%
  \setlength{\itemsep}{0pt}\setlength{\parskip}{0pt}}
\setcounter{secnumdepth}{5}
\ifLuaTeX
  \usepackage{selnolig}  % disable illegal ligatures
\fi
\IfFileExists{bookmark.sty}{\usepackage{bookmark}}{\usepackage{hyperref}}
\IfFileExists{xurl.sty}{\usepackage{xurl}}{} % add URL line breaks if available
\urlstyle{same}
\hypersetup{
  colorlinks=true,
  linkcolor={blue},
  filecolor={Maroon},
  citecolor={Blue},
  urlcolor={blue},
  pdfcreator={LaTeX via pandoc}}

\author{}
\date{}

\begin{document}

\hypertarget{chapter-07-ux884cux5217ux5206ux89e3}{%
\section{Chapter 07 ---
行列分解}\label{chapter-07-ux884cux5217ux5206ux89e3}}

この章は前半の概念が合流する場所です。

\hypertarget{ux306aux305cux5206ux89e3ux3059ux308bux306eux304b}{%
\subsection{なぜ分解するのか}\label{ux306aux305cux5206ux89e3ux3059ux308bux306eux304b}}

一つの大きな操作 \texttt{A} を、そのまま理解するより

\[
A=A_1A_2\cdots A_k
\]

と意味の分かる操作に分ける方が理解しやすく、計算にも有利です。

\hypertarget{lu-qr-evd-svd-ux306eux5f79ux5272}{%
\subsection{LU / QR / EVD / SVD
の役割}\label{lu-qr-evd-svd-ux306eux5f79ux5272}}

\hypertarget{lu}{%
\subsubsection{LU}\label{lu}}

Gauss 消去を保存した分解。主目的は連立方程式の効率的な解法です。

\hypertarget{qr}{%
\subsubsection{QR}\label{qr}}

列空間を正規直交基底 \texttt{Q} と係数 \texttt{R}
に分けます。最小二乗に強い分解です。

\hypertarget{evd}{%
\subsubsection{EVD}\label{evd}}

同じ空間内の不変方向を探します。反復変換や力学系を理解するのに適します。

\hypertarget{svd}{%
\subsubsection{SVD}\label{svd}}

\[
A=U\Sigma V^T
\]

任意の行列を、入力側の直交基底・純粋な伸縮・出力側の直交基底へ分けます。

\hypertarget{svd-ux3092ux7406ux89e3ux3059ux308bux6838ux5fc3}{%
\subsection{SVD
を理解する核心}\label{svd-ux3092ux7406ux89e3ux3059ux308bux6838ux5fc3}}

必ず次の式を言葉にしてください。

\[
Av_i=\sigma_i u_i.
\]

これは「入力方向 \texttt{v\_i} を入れると、方向 \texttt{u\_i}
へ向かう長さ \texttt{σ\_i} の出力になる」という意味です。

さらに

\[
A=\sum_i\sigma_i u_iv_i^T
\]

と書けば、行列は重要度順の rank-1 成分へ分解されます。

ここから低ランク近似が出ます。\texttt{σ\_i=0}
の方向は核になり、\texttt{1/σ\_i}
を使えば擬似逆が出ます。\texttt{σ\_max/σ\_min}
を見れば条件数になります。

\hypertarget{ux5230ux9054ux57faux6e96}{%
\subsection{到達基準}\label{ux5230ux9054ux57faux6e96}}

SVD から rank、null space、pseudoinverse、condition number、rank-k
approximation を一続きで導けることを目標にします。

\end{document}
